\section{One source, two continuers}
\label{sec:branching}

The fission cases of personal-identity theory ask what happens when one source
divides into two equally good successors \citep{Parfit1984, Olson2023SEP}. We
cannot settle that question and do not try. We can say exactly what the
structural checks of this paper do in such a case. The short answer is that
continuation becomes a relation with two continuers, and nothing in the
structure picks one of them.

\paragraph{A symmetric fork that passes.}
The chain model of Section~\ref{sec:continuation} also has a fork. One source
system hands off to two target systems, and both branches receive identical
inputs, so the fork is symmetric by construction.

\begin{theorem}[Continuation is relation-valued]
% claim: TTH-08
\label{thm:fork}
In the symmetric fork of the chain model, labelled
\snapGSevenForkSystemIds{} (source first), the following hold.
\begin{enumerate}
\item The source continues to each of the two targets, and the valid targets
are exactly these two. Each carries the chain model's handoff certificate,
covering realization, lineage and autonomy after withdrawal, for every state of
the finite factor.
\item The swap of the two targets preserves validity and is an involution.
\item No continuer is fixed by the swap.
\end{enumerate}
\end{theorem}

The substance of the theorem is the first part. Each branch carries the chain
model's handoff certificate: its realization check (both of its closed condition checks
accepted, faithful receiver content, an actual entry at the target's root), its
lineage check (future-transport laws, a passing handoff, raw transport), and a
controller that wins its game after withdrawal. This is a kernel-checked
computation on the model. It is the model's own certificate, not the full
generic gate of Definition~\ref{def:certificate}. The third part follows at once, since the swap exchanges the only
two continuers. Its consequence is the one that matters. No rule that commutes
with the swap can select a single continuer, because the swap fixes neither.
This does not make either choice invalid: choosing one target alone still
gives a valid continuation. What fails is any ground in the structure for the
choice.

A smaller fork shows the same thing with less machinery.

\begin{proposition}[Symmetry gives no ground for choosing]
% claim: TTN-11
\label{prop:symmetric-fork}
On the two-copy model there is a symmetric fork into two distinct target
systems, both passing the realization and autonomy checks, in which the swap
maps valid continuers to valid continuers and fixes none of them.
\end{proposition}

Again the argument against a symmetric choice is elementary: the swap has no
fixed point at all. We include the proposition because it shows that symmetric
forks already arise for models far simpler than the chain.

\paragraph{Ancestry does not decide identity.}
One might hope to recover a single continuer from ancestry, that is, from
which target descends from which source. On the chain model, after the fork,
we can test this directly, with ancestry computed from the actual execution as
in Section~\ref{sec:ancestry}.

\begin{theorem}[Ancestry underdetermines identity]
% claim: TTF-07
\label{thm:identity}
On the chain model, for every $n$, there is a safe execution of $n+1$ steps
after withdrawal, ending in a fission, in which the source has two descendants
holding the same state. That state passes the model's realization
check, with conditions (4) to (7) in simplified form, and obeys its
continuation law.
Consider identity relations that are equivalences, that identify the source
with at most one other node, and that identify the source with no node other
than itself that is not its descendant under the computed ancestry. Among these there are two that
identify the source with different descendants, and neither identification is
forced.
\end{theorem}

The identity notion here is deliberately weak and explicit: an equivalence
relation that pairs the source with itself and with at most one other node,
which must be a descendant. The theorem says that ancestry, even computed from an actual
execution, does not choose between the two descendants. It does not say that
neither is the source, nor that identity is undefined. A stronger statement,
that ancestry rules out identity under every possible notion, is false as
stated, and we do not make it.

These results fit the main positions in the literature without deciding
between them. On the closest-continuer theory a tie leaves no continuer at all
\citep{Nozick1981}, which is consistent with our result that no symmetric rule
picks one. Such a theory makes whether one candidate continues the source
depend on whether a rival exists. That is the standard criticism of it, and
\citet{Hawley2005} examines what exactly is wrong with it. On the multiple-occupancy view two persons share their stages
before the division \citep{Lewis1976}, and we have instead let the
continuation relation have two values. \citet{Cerullo2015} argues informally
that uploading yields branching continuers with no reason to favour either. The
theorems above are formal structural analogues of that view on finite systems.
Our identity notion, however, is constrained by ancestry and is shown not to
decide, whereas his criterion is built from ancestry and is taken to settle
continuity. \citet{Weber2025Multiplicity} states the same structure for
uploading: the continuity relations that underlie persistence can hold from
one to many. He argues that it is plausible that we could become several distinct
persons, and we take no view on that. Nothing here is a claim about persons.

\paragraph{Branches preserved, recombination lost.}
Fission raises a converse question: if each branch is preserved, is the way the
branches combine preserved too? It need not be.

\begin{proposition}[Branchwise preservation is not recombination]
% claim: TTN-14
\label{prop:recombination}
On the two-copy model, the staged map preserves every ordered branch record
from source to target, yet two things fail. On the target, the two assemblages
have equal joint couplings, so joint content no longer strictly refines the
branch records. And the source's and the target's joint couplings differ, so
the staged transport of the binding condition fails.
\end{proposition}

The binding condition here is in its simplified form, a declared coupling, so
the proposition concerns that form. It is the cross-system counterpart of a
familiar point: aggregating branch contents does not determine their
combination \citep{TsiokosArchitecture2026}. Figure~\ref{fig:fork} draws the
fork and the swap.

\begin{figure}[tbp]
\centering
% F6: the symmetry exchanges targets and applying it twice returns each target.
\begin{tikzpicture}[x=1cm,y=1cm,font=\small,
 system/.style={draw,rounded corners=2pt,fill=black!3,minimum width=3.1cm,minimum height=.95cm,align=center},
 flow/.style={-{Stealth[length=2mm]},semithick}]
\node[system] (source) at (1.8,1.2) {Source system\\\snapGSevenForkSystemIdsItemZero};
\node[system] (left) at (8,3.1) {Target system\\\snapGSevenForkSystemIdsItemOne};
\node[system] (right) at (8,-.7) {Target system\\\snapGSevenForkSystemIdsItemTwo};
\draw[flow] (source.east) to[out=0,in=180] node[above,sloped] {handoff} (left.west);
\draw[flow] (source.east) to[out=0,in=180] node[below,sloped] {handoff} (right.west);
\draw[<->,densely dashed,line width=1pt] (left.east) to[out=0,in=0]
 node[right,align=center] {swap\\involution} (right.east);
\node[align=center,text width=3.5cm] at (13,3.1) {autonomous\\after withdrawal};
\node[align=center,text width=3.5cm] at (13,-.7) {autonomous\\after withdrawal};
\node[draw,densely dotted,align=center,text width=3.1cm] at (1.8,-.7) {Source and transfer\\aids withdrawn};
\node[align=center] at (8,-1.9) {Swap twice: each target returns to itself\\No branch is structurally selected as the original};
\end{tikzpicture}

\caption{The symmetric fork. One source hands off to two targets that receive identical inputs, and each is autonomous after withdrawal. The swap of the two targets is an involution and fixes neither, so nothing in the structure selects one of them. We draw no conclusion about the identity of either branch.}
\label{fig:fork}
\end{figure}
